\section{Odd inertia and the full CM extension obstruction}
\label{sec:general-odd-inertia}

The scalar coordinates admit exact product descriptions even when
their individual full images are unknown. An inertia argument then
eliminates the conjugation obstruction for the entire relative
character module. The relative central extension class remains
the condition for extending an intrinsic character.

Put \(R_0=R_{\rm quad}\), \(M=\ker\lambda\), and \(D=N_{\rm rad}\).
The definitions in
Theorems~\ref{thm:simultaneous-quadratic-normalization}
and~\ref{thm:complete-quadratic-normalization} give
\begin{equation}
 \begin{split}
 R_0={}&D\cap\ker\Phi_2\cap\ker\mathfrak t_2\\
 &\cap\bigcap_{\ell<q,\ \ell,q\ {\rm odd}}\ker\xi_{\ell,q}
 \cap\bigcap_{p\ {\rm odd},\ p\ne3}\ker\nu_p,
 \qquad G=R_0H.
 \end{split}
 \label{eq:general-cm-residual-intersection}
\end{equation}
Here \(\Phi_2\) is the full based free-pro-two action and
\(\mathfrak t_2\) is the entire corrected cubic coefficient-two
homomorphism.

For \(q\ge5\), write \(\rho^{\rm bas}_{q,2}\) for the full
compact-homology representation of the degree-\(q\) cyclic cover,
and \(\Delta_q\) for its deck generator.
Equation~\eqref{eq:coefficient-two-cross-3} gives
\[
 \rho^{\rm bas}_{q,2}(g)\Delta_q
       \rho^{\rm bas}_{q,2}(g)^{-1}
       =\Delta_q^{\lambda(g)\bmod q}.
\]
The integral centralizer of \(\Delta_q\) is the commutative
algebra \(\mathcal O_q=\mathbb Z[\zeta_q]\otimes\mathbb Z_2\),
since every nontrivial deck eigencharacter occurs once.
On \(M\), the representation is scalar in this algebra.
Its principal pro-two projection is \(h_q\).
Projection to principal units commutes with the cyclotomic
automorphisms, including permutations of the local factors.
Consequently
\begin{equation}
 h_q(gmg^{-1})=[\lambda(g)\bmod q]\bigl(h_q(m)\bigr)
                              \qquad(g\in G,\ m\in M).
 \label{eq:general-cm-scalar-equivariance}
\end{equation}
In particular every \(\ker h_q\) is closed and normal in \(G\).

\subsection{A prime outside the imposed scalar kernels}

For a set \(\Lambda\) of primes at least five define
\[
 R_\Lambda=R_0\cap\bigcap_{q\in\Lambda}\ker h_q.
\]
The next theorem assumes arithmetic factorization for this
particular intersection.

\begin{theorem}\label{thm:omitted-prime-vanishing}
Suppose \(G=R_\Lambda H\), and let \(p\ge5\) be a prime
outside \(\Lambda\).
Let \(W\) be a finite abelian \(2\)-primary
\(G/R_\Lambda\)-module whose action factors through
\[
 G/R_\Lambda\xrightarrow{\lambda\bmod p}
                  (\mathbb Z/p)^\times=\langle s_p\rangle.
\]
If \(s_p-1\) is bijective on \(W\), then
\begin{equation}
                  H^j_{\rm cont}(G/R_\Lambda,W)=0
                                     \qquad(j\ge0).
 \label{eq:omitted-prime-cohomology}
\end{equation}
\end{theorem}

\begin{proof}
We apply Proposition~\ref{prop:odd-inertia-vanishing}.
Define
\begin{equation}
 L_\Lambda=D\cap\ker\Phi_2\cap\ker\mathfrak t_2
                    \cap\bigcap_{q\in\Lambda}\ker h_q.
 \label{eq:general-cm-intermediate-subgroup}
\end{equation}
Every kernel here is closed and \(G\)-normal.
For \(D\), use the full affine Kummer homomorphism of
Theorem~\ref{thm:rational-radical-normalization}.
The map \(\Phi_2\) is a representation on \(G\).
For the corrected cubic frame, the integral conjugation
identity \eqref{eq:adelic-cubic-kernel-conjugation}
preserves its identity fibre in \(M\).
The remaining kernels are normal by
\eqref{eq:general-cm-scalar-equivariance}.

The two quadratic families in
\eqref{eq:general-cm-residual-intersection} have joint
kernel \(R_\Lambda\) on \(L_\Lambda\).
Their invariance follows from
Theorems~\ref{thm:all-odd-cross-cup-lifts}
and~\ref{thm:coefficient-two-cross-cups}.
Thus \(L_\Lambda/R_\Lambda\) is a closed subgroup of
a product of copies of \(\mathbb F_2\), and
\([G,L_\Lambda]\subseteq R_\Lambda\).
In particular it is central pro-two in \(G/R_\Lambda\).
The intersection \eqref{eq:general-cm-intermediate-subgroup}
includes the kernel of the entire cubic coefficient-two map.

\paragraph{The odd part of arithmetic inertia.}
Write \(p-1=2^e m\), with \(m\) odd.
Let \(Z_p\subset\widehat{\mathbb Z}^{\times}\) be the
finite order-\(m\) subgroup in the \(p\)-component,
with all other components equal to \(1\).
Cyclotomic inertia at \(p\) maps onto
\(\mathbb Z_p^\times\) and is trivial at every other
coefficient prime: the \(p\)-power cyclotomic tower
is totally ramified locally, whereas roots of unity
of order prime to \(p\) are unramified.
Choose \(b\in I_p\subset H\) whose multiplier generates
\(Z_p\).
Let \(e_{\rm odd}\in\widehat{\mathbb Z}\) be zero at
two and one at every odd prime, and put
\(a=b^{e_{\rm odd}}\), with powers taken in the
procyclic closure of \(b\).
Then
\[
 \overline{\langle a\rangle}\text{ is pro-odd},\qquad
 \lambda(a)\text{ generates }Z_p,\qquad
 \lambda_\ell(a)=1\quad(\ell\ne p).
\]

We have \(W^a=0\).
Indeed \(a\) acts through a generator of the odd-order
subgroup of \(\langle s_p\rangle\).
If \(V=W^{\langle s_p^{2^e}\rangle}\) were nonzero,
then \(V[2]\) would be a nonzero finite
\(\mathbb F_2\)-vector space.
The action of \(s_p\) on it has order dividing \(2^e\),
so \((s_p-1)^{2^e}=0\).
This contradicts the injectivity of \(s_p-1\).
Thus \(V=0\), including the case \(m=1\), when the
hypothesis forces \(W=0\).

Let \(\mathscr N_H\) be the closed normal closure of
\(a\) in \(H\).
We claim
\begin{equation}
                       \mathscr N_H\cap D\subseteq L_\Lambda.
 \label{eq:general-cm-arithmetic-kernel-comparison}
\end{equation}

\paragraph{The based action and the corrected cubic frame.}
The full based image \(\Phi_2(G)\) is pro-two.
The peripheral formulas
\(x\mapsto x^\lambda\), \(y\mapsto f^{-1}y^\lambda f\)
act trivially on the mod-two Frattini quotient of
the free pro-two group, because \(\lambda\) is odd.
The kernel of the corresponding automorphism action
is pro-two, as seen on finite characteristic
two-group quotients; this is the argument in the
proof of Theorem~\ref{thm:full-dyadic-free-product}.
It follows that \(\Phi_2(a)=1\), and hence
\(\Phi_2(\mathscr N_H)=1\).
This is an equality of full based automorphisms.

The curve \(E_3:Y^2=X^3+1\) has discriminant \(-432\),
so it has good reduction at \(p\ge5\).
Its two-adic Tate representation is unramified there.
The Galois field
\(\mathbb Q(\mu_4,3^{1/4})\) is also unramified at \(p\):
the discriminant of \(X^4-3\) has prime divisors only
\(2,3\), and adjoining \(\mu_4\) adds ramification
only at two.
Both arithmetic kernels contain \(\mathscr N_H\).

For \(\sigma\in\mathscr N_H\cap M\), the signed comparison
of Proposition~\ref{prop:adelic-signed-path-comparison}
therefore gives
\[
 \mathfrak t_2(\sigma)
   =\frac{\sigma(3^{1/4})}{3^{1/4}}\rho_{E_3,2}(\sigma)=1.
\]
This comparison is applied only on its domain \(H\cap M\),
with its quartic Kummer correction included.

\paragraph{Good reduction and the odd deck factors.}
Fix \(q\in\Lambda\), so \(q\ne p\).
The curve \(w^q=t(t-1)\) has good reduction at \(p\).
For a direct check, write \(q=2g+1\), set \(y=2t-1\),
and use the charts
\[
 y^2=1+4w^q,\qquad
 v^2=u^{q+1}+4u,\qquad
 u=w^{-1},\quad v=y/w^{g+1}.
\]
The first branch polynomial is separable modulo \(p\),
since \(p\nmid2q\).
At the point \(u=v=0\) of the second chart the derivative
with respect to \(u\) is \(-4\), a unit.
The two monic quadratic charts give a finite flat cover
of \(\mathbb P^1_{\mathbb Z_p}\) which is proper and smooth.
Its Jacobian is an abelian scheme.
Its \(2^n\)-torsion is finite etale for every \(n\), hence
constant over the strict henselization.
This proves unramifiedness of the entire two-adic
Tate representation, both here and for \(E_3\).

The arithmetic Tate representation and the based
compact-homology representation differ only by a
deck element, as in the proof of
Proposition~\ref{prop:pdl-arithmetic-realization}.
Indeed two lifts of the same outer arithmetic action
differ by an inner automorphism of the base geometric
group. The covering-subgroup part acts trivially on
compact homology, leaving its deck class.
Since inertia at \(p\) fixes \(\mathbb Q(\zeta_q)\),
good reduction gives
\[
                 \rho^{\rm bas}_{q,2}(a)\in
                              \langle\Delta_q\rangle\simeq C_q.
\]
This deck element need not be trivial.
The full normalizer identity shows that
\(\langle\Delta_q\rangle\) is normalized by
\(\rho^{\rm bas}_{q,2}(H)\).
Consequently
\[
                 \rho^{\rm bas}_{q,2}(\mathscr N_H)
                                  \subseteq\langle\Delta_q\rangle.
\]
For \(\sigma\in\mathscr N_H\cap M\), its principal
pro-two scalar \(h_q(\sigma)\) is therefore one,
because that projection kills the odd-order deck group.
This holds for every imposed coordinate, including
an infinite collection.
Together with \(D\subseteq M\), it proves
\eqref{eq:general-cm-arithmetic-kernel-comparison}.

\paragraph{The full affine radical quotient.}
By Theorem~\ref{thm:rational-radical-normalization},
\[
 G/D\simeq A_{\rm rad}
  =\left\{(u,(t_r)_r):
      u\in\widehat{\mathbb Z}^{\times},\
      t_r\in\widehat{\mathbb Z},\
      t_r\bmod2=\epsilon_r(u)\right\},
 \qquad (u,t)(v,z)=(uv,t+uz),
\]
where \(r\) ranges over all rational primes and
\(\epsilon_r\) is the character of \(\mathbb Q(\sqrt r)\).
The subgroup
\[
 \mathscr O_p
   =\{(u,t)\in A_{\rm rad}:u\in Z_p,\
                    t_{r,2}=0\text{ for every }r\}
\]
is closed and pro-odd.
Every \(\epsilon_r\) vanishes on \(Z_p\), so it is
the semidirect product
\[
 \left(\prod_{r\ {\rm prime}}\prod_{\ell\ {\rm odd}}\mathbb Z_\ell\right)
                                      \rtimes Z_p.
\]
It is normal in the entire affine group.
Indeed conjugation by \((v,z)\) sends \((u,t)\) to
\[
                            (u,vt+(1-u)z);
\]
every dyadic translation stays zero because \(u_2=1\).

The image of \(a\) belongs to \(\mathscr O_p\).
Its multiplier belongs to \(Z_p\), and every
dyadic Kummer cocycle restricts on
\(\overline{\langle a\rangle}\) to a homomorphism
to the additive group \(\mathbb Z_2\).
That homomorphism is zero by coprimality.
The full image of \(\mathscr N_H\) is therefore
contained in \(\mathscr O_p\), and
\[
                        \mathscr N_H/(\mathscr N_H\cap D)
                                         \text{ is pro-odd}.
\]
This argument includes all rational-prime translations
and all odd coefficient components.

Finally the assumed factorization \(G=R_\Lambda H\)
implies \(G=L_\Lambda H\), and compactness gives
\[
                      G/L_\Lambda\simeq H/(H\cap L_\Lambda).
\]
The closed normal closure of \(a\) in this quotient
is exactly the image of \(\mathscr N_H\).
Equation~\eqref{eq:general-cm-arithmetic-kernel-comparison}
and \(L_\Lambda\subseteq D\) give
\[
                   \mathscr N_H\cap L_\Lambda=\mathscr N_H\cap D.
\]
Its normal closure in \(G/L_\Lambda\) is therefore
pro-odd. We have verified the central pro-two layer,
the pro-odd procyclic closure of \(a\), its pro-odd
normal closure in \(G/L_\Lambda\), and \(W^a=0\).
Proposition~\ref{prop:odd-inertia-vanishing} proves
\eqref{eq:omitted-prime-cohomology}.
\end{proof}

\subsection{Exact products of the scalar images}

Let \(\mathcal I\) be any set of primes at least five.
Define the full and arithmetic images by
\begin{equation}
 \begin{aligned}
 C_p&=h_p(R_0),&
 D_p&=h_p(R_0\cap H),\\
 C&=(h_p)_{p\in\mathcal I}(R_0),&
 D_0&=(h_p)_{p\in\mathcal I}(R_0\cap H).
 \end{aligned}
 \label{eq:all-cm-scalar-images}
\end{equation}
All these groups are compact abelian pro-two groups.
The group \(C\) is \(G\)-stable.
The group \(D_0\) is \(H\)-stable and is also \(G\)-stable:
conjugation by \(R_0\) is trivial on \(C\), and \(G=R_0H\).
Their conjugation actions therefore factor through
\(\Gamma_0=G/R_0\).

Use additive operator notation on the scalar unit groups.
For each \(p\), retain
\[
 \mathcal O_p=\mathbb Z[\zeta_p]\otimes\mathbb Z_2,\qquad
 \mathcal U_p=1+2\mathcal O_p,\qquad
 T_p=\{t\in\mathcal U_p:t\,c(t)=1\},\qquad P_p=(s_p-1)T_p.
\]
Here \(s_p\) is a primitive cyclotomic generator and
\(c=s_p^{(p-1)/2}\).
Proposition~\ref{prop:explicit-cm-type-factor} and
Theorem~\ref{thm:cm-type-torus-unit} give the exact identities
\begin{equation}
 (s_p-1)\Phi_p^*=(c-1)Q_p,\qquad
 (c-1)\mathcal U_p=T_p,\qquad
 Q_p(1)=\varepsilon_p\in\{1,-1\}.
 \label{eq:all-cm-integral-operator-inputs}
\end{equation}
The upper bound \eqref{eq:pdl-upper-input} gives
\(C_p\subseteq P_p\).

\begin{lemma}\label{lem:all-cm-isolated-axes}
If \(i_p:P_p\to\prod_{q\in\mathcal I}P_q\)
denotes inclusion in the \(p\)-coordinate, then
\[
                           i_p(Q_pP_p)\subseteq D_0\subseteq C.
\]
In particular \(Q_pP_p\subseteq D_p\) and \(Q_pC_p\subseteq D_p\).
\end{lemma}

\begin{proof}
The witnesses of Proposition~\ref{prop:pdl-arithmetic-realization}
lie in \(V_{\rm quad}\cap H\subseteq R_0\cap H\),
have every other prime coordinate equal to one, and realize
\((s_p-1)^2\Phi_p^*(\mathcal U_p)\) in degree \(p\).
By \eqref{eq:all-cm-integral-operator-inputs}, this image is
\[
 (s_p-1)(c-1)Q_p(\mathcal U_p)
     =Q_p(s_p-1)T_p=Q_pP_p.
\]
These are equalities of the full compact groups,
including their finite sign subgroups.
\end{proof}

\begin{proposition}\label{prop:all-cm-exact-products}
For every set \(\mathcal I\) of primes at least five,
\begin{equation}
                         C=\prod_{p\in\mathcal I}C_p,\qquad
                         D_0=\prod_{p\in\mathcal I}D_p.
 \label{eq:all-cm-exact-products}
\end{equation}
Consequently
\[
                         C/D_0\simeq\prod_{p\in\mathcal I}C_p/D_p.
\]
These statements require no CM-unit or finite-index hypothesis.
\end{proposition}

\begin{proof}
Choose an arithmetic inertia element at \(p\) whose
cyclotomic reduction modulo \(p\) is primitive.
Its induced operator \(S_p\) on the ambient scalar product
acts by \(s_p\) in the \(p\)-coordinate and by the identity
in every other coordinate, since cyclotomic inertia at \(p\)
is trivial modulo every \(q\ne p\).
Both \(C\) and \(D_0\) are stable under \(S_p\).

Divide the integral polynomial \(Q_p(X)-\varepsilon_p\)
by \(X-1\), and write
\begin{equation}
                     Q_p(X)=\varepsilon_p+(X-1)F_p(X),
                     \qquad F_p\in\mathbb Z[X].
 \label{eq:all-cm-polynomial-division}
\end{equation}
For \(x\in C\), the following identity holds in the whole
ambient product:
\begin{equation}
 \varepsilon_p i_p(x_p)
       =i_p(Q_px_p)-(S_p-1)F_p(S_p)x.
 \label{eq:all-cm-exact-coordinate-projection}
\end{equation}
At \(p\), it is \eqref{eq:all-cm-polynomial-division};
at every other coordinate the second term is zero.
The first term belongs to \(D_0\) by
Lemma~\ref{lem:all-cm-isolated-axes}, and the second
belongs to \(C\) by stability.
Since \(\varepsilon_p=\pm1\), it follows that
\(i_p(C_p)\subseteq C\).

For \(x\in D_0\), both terms on the right of
\eqref{eq:all-cm-exact-coordinate-projection}
belong to \(D_0\), so \(i_p(D_p)\subseteq D_0\).
Finite sums of coordinate axes are dense in the
corresponding products. The images \(C,D_0\) are
compact and hence closed, proving
\eqref{eq:all-cm-exact-products}.
This also covers infinite and empty index sets.
The quotient isomorphism is induced by the coordinate
quotient maps and is a continuous bijection between
compact Hausdorff groups.
\end{proof}

The proposition describes the full and arithmetic images \(C_p,D_p\).
It does not assert that either of them equals \(P_p\).

\subsection{The full continuous conjugation module}

Set
\[
 E_p=C_p/D_p,\qquad
 U_p=\operatorname{Hom}_{\rm cont}(E_p,\mathbb F_2),\qquad
 U=\operatorname{Hom}_{\rm cont}(C/D_0,\mathbb F_2).
\]
Give the character groups the discrete topology and the
left conjugation action
\((g\chi)(x)=\chi(g^{-1}xg)\).
Proposition~\ref{prop:all-cm-exact-products} gives
\begin{equation}
 C/D_0\simeq\prod_{p\in\mathcal I}E_p,\qquad
 U\simeq\bigoplus_{p\in\mathcal I}U_p.
 \label{eq:all-cm-character-decomposition}
\end{equation}
Indeed the open kernel of a continuous character on the
product contains the entire factors outside a finite set.
Conversely every finite sum of coordinate characters is
continuous. The decomposition is equivariant, and the
action on \(U_p\) factors through \(\lambda\bmod p\).

\begin{lemma}\label{lem:all-cm-coordinate-module}
Each \(U_p\) is finite, and \(s_p-1\) is bijective on it.
\end{lemma}

\begin{proof}
The principal-unit group \(1+2\mathcal O_p\) has the
finite-index subgroup \(1+4\mathcal O_p\), which logarithm
identifies with \(4\mathcal O_p\simeq\mathbb Z_2^{p-1}\).
It is therefore a finitely generated abelian pro-two group.
A closed subgroup has the same finiteness property:
its intersection with the logarithmic subgroup is a
closed submodule of a finite free \(\mathbb Z_2\)-module,
and its remaining quotient is finite.
Thus \(C_p\) and its quotient \(E_p\) are finitely generated,
so \(U_p\) is finite. This argument retains all finite torsion.

By Lemma~\ref{lem:all-cm-isolated-axes}, \(Q_p\) annihilates
\(E_p\). Equation~\eqref{eq:all-cm-polynomial-division}
therefore gives
\[
 (s_p-1)(-\varepsilon_p F_p(s_p))
       =1=(-\varepsilon_p F_p(s_p))(s_p-1)
                         \quad\text{on }E_p.
\]
Hence \(s_p-1\) is a continuous automorphism of \(E_p\).
The dual action is precomposition with \(s_p^{-1}\),
and \(s_p^{-1}-1=-s_p^{-1}(s_p-1)\) is also invertible.
Duality proves the assertion on \(U_p\).
\end{proof}

\begin{corollary}\label{cor:all-cm-conjugation-vanishing}
With its discrete conjugation action, \(U\) is a continuous
\(\Gamma_0=G/R_0\)-module, and
\begin{equation}
                       H^j_{\rm cont}(\Gamma_0,U)=0
                                      \qquad(j\ge0).
 \label{eq:all-cm-conjugation-cohomology}
\end{equation}
\end{corollary}

\begin{proof}
Each element of \(U\) has finite coordinate support.
The action on each coordinate factors through a finite
cyclotomic group, so each element has an open stabilizer.
This proves continuity for the discrete topology.

Apply Theorem~\ref{thm:omitted-prime-vanishing} with
\(\Lambda=\varnothing\) and \(R_\Lambda=R_0\).
The required factorization is
\eqref{eq:general-cm-residual-intersection}, and
Lemma~\ref{lem:all-cm-coordinate-module} supplies the
coefficient hypotheses. Thus
\(H^j(\Gamma_0,U_p)=0\) for every \(p,j\).
The same holds for every finite sum
\(U_F=\bigoplus_{p\in F}U_p\).

For every \(n\ge0\), a continuous cochain
\(\Gamma_0^n\to U\) has finite image because its domain
is compact and \(U\) is discrete.
The union of the supports of those finitely many values
is finite. Consequently, as cochain complexes,
\begin{equation}
 C^\bullet_{\rm cont}(\Gamma_0,U)
      =\bigcup_{\substack{F\subseteq\mathcal I\\F\ {\rm finite}}}
                            C^\bullet_{\rm cont}(\Gamma_0,U_F).
 \label{eq:all-cm-continuous-cochain-union}
\end{equation}
Every differential preserves these subcomplexes.
A positive-degree cocycle therefore belongs to a finite
summand complex and is a coboundary there.
In degree zero an invariant element likewise belongs
to a finite summand with zero invariant space.
This proves \eqref{eq:all-cm-conjugation-cohomology}
for the full, possibly infinite module.
\end{proof}

\subsection{The remaining relative central extension class}

Let
\[
 J=\ker\bigl((h_p)_{p\in\mathcal I}:R_0\to C\bigr),\qquad
 H_0=R_0\cap H,
\]
and let
\[
 \mathscr X(J)=
 \{\kappa:J\to\mathbb F_2:
   \kappa\text{ is continuous and }G\text{-invariant},\
   \kappa(J\cap H)=0\}.
\]
For \(\mathcal I\) equal to all primes at least five,
one has \(J=R_{\rm all}\).

\begin{theorem}\label{thm:all-cm-central-extension}
For \(\kappa\in\mathscr X(J)\), the central extension
\(R_0/\ker\kappa\to C\), with its prescribed arithmetic
splitting over \(D_0\), defines a continuous relative class
\[
                         o_2(\kappa)\in
                                H^2_{\rm cont}(C,D_0;\mathbb F_2)
\]
when \(\kappa\ne0\); set \(o_2(0)=0\).
The assignment \(o_2\) is \(\mathbb F_2\)-linear.
The character \(\kappa\) extends to a continuous character
of \(G\) vanishing on \(H\) if and only if \(o_2(\kappa)=0\).
When it exists, this extension is unique.
Equivalently, restriction identifies
\begin{equation}
 \operatorname{Hom}_{\rm cont}(B,\mathbb F_2)
   \simeq
 \ker\bigl(o_2:\mathscr X(J)\longrightarrow
                       H^2_{\rm cont}(C,D_0;\mathbb F_2)\bigr).
 \label{eq:all-cm-central-kernel-identification}
\end{equation}
\end{theorem}

\begin{proof}
For \(\kappa\ne0\), \(G\)-invariance makes \(\ker\kappa\)
normal in \(G\) and gives a central extension
\[
 1\longrightarrow\mathbb F_2
   \longrightarrow E:=R_0/\ker\kappa
   \xrightarrow{\pi}C\longrightarrow1.
\]
The image of \(H_0\) in \(E\) maps isomorphically onto
\(D_0\): its kernel comes from \(J\cap H\), where
\(\kappa\) is zero.
Compactness gives a prescribed continuous splitting
\(\sigma_D:D_0\to E\).

We make the continuous relative class explicit.
A finite-kernel profinite extension admits a continuous
set-theoretic section.
Choose an open normal subgroup of \(E\) meeting the finite
kernel trivially. It maps homeomorphically onto its open
image in \(C\), and finitely many translated copies give
a section on all of \(C\).
The discrepancy between this section on \(D_0\) and
\(\sigma_D\) is a continuous finite-valued function.
A finite clopen partition of a closed subset of a
profinite space extends to a finite clopen partition
of the whole space, by disjoint clopen separation.
Thus the discrepancy extends to \(C\), and adjusting
the section gives a continuous section \(s:C\to E\)
agreeing with \(\sigma_D\).

Its continuous cocycle
\(s(x)s(y)s(xy)^{-1}\) has values in \(\mathbb F_2\)
and vanishes on \(D_0\times D_0\).
Changing \(s\) while keeping this prescribed restriction
changes the cocycle by the coboundary of a continuous
one-cochain zero on \(D_0\).
This describes its relative class.
It also agrees with the relative mapping-cone definition:
restriction of continuous finite-valued cochains from
\(C^n\) to \(D_0^n\) is surjective by the same clopen
extension argument, so the kernel relative complex
computes the relative cohomology.
The class is zero exactly when a continuous
one-cochain zero on \(D_0\) changes \(s\) into a
homomorphic splitting extending \(\sigma_D\).

A compatible splitting identifies
\(E\simeq\mathbb F_2\times C\).
Projection onto the first factor, pulled back to \(R_0\),
gives a continuous character \(\ell:R_0\to\mathbb F_2\)
extending \(\kappa\) and zero on \(H_0\).
Conversely such a character factors through \(E\),
is the identity on the central \(\mathbb F_2\),
and kills \(\sigma_D(D_0)\).
Its kernel maps homeomorphically onto \(C\), giving
the required compatible splitting.
For \(\kappa=0\), the zero character provides such
an extension and \(o_2(0)=0\).
Pushout and Baer sum preserve the prescribed arithmetic
splitting. They identify the class associated with
\(\kappa+\kappa'\) with \(o_2(\kappa)+o_2(\kappa')\),
so \(o_2\) is \(\mathbb F_2\)-linear.

Suppose \(o_2(\kappa)=0\).
The possible characters \(\ell\) on \(R_0\) form a torsor
under the full module
\(U=\operatorname{Hom}_{\rm cont}(C/D_0,\mathbb F_2)\):
the difference of two extensions kills \(J\), factors
continuously through \(C\), and kills \(D_0\).
Conjugation by \(G\) preserves this torsor.
The subgroup \(H\) preserves \(H_0\) and the specified
restriction on \(J\), while inner conjugation by \(R_0\)
acts trivially on characters of \(R_0\).
Since \(G=R_0H\), the action factors through \(\Gamma_0\).

This is a continuous action on a discrete torsor, even
when \(U\) is infinite.
For any continuous \(\ell:R_0\to\mathbb F_2\), choose an
open normal subgroup \(V\triangleleft G\) with
\(V\cap R_0\subseteq\ker\ell\).
Then \([V,R_0]\subseteq V\cap R_0\), so \(V\) fixes \(\ell\).
The stabilizer, and its image in \(\Gamma_0\), is open.

Choose one extension \(\ell_0\).
The differences
\[
                         z(\gamma)=\gamma\ell_0-\ell_0
                                      \qquad(\gamma\in\Gamma_0)
\]
form a continuous \(U\)-valued one-cocycle.
Corollary~\ref{cor:all-cm-conjugation-vanishing} gives
\(H^1(\Gamma_0,U)=0\), so
\(z(\gamma)=\gamma u-u\) for some \(u\in U\).
The extension \(\ell_0-u\) is \(G\)-invariant.
It is the unique invariant extension, because the same
corollary gives \(U^{\Gamma_0}=0\).

Finally define
\[
                         \chi(rh)=\ell(r)
                            \qquad(r\in R_0,\ h\in H)
\]
using this invariant extension.
The condition \(\ell(H_0)=0\) proves well-definedness;
\(H\)-invariance proves the homomorphism identity.
Its pullback along the compact surjection
\(R_0\times H\to G\) is continuous, so \(\chi\) is continuous.
It extends \(\kappa\) and vanishes on \(H\).
Conversely a character of \(G\) with these properties
restricts to a compatible character on \(R_0\), forcing
\(o_2(\kappa)=0\).
The difference of two such characters is zero on \(R_0\)
by \(U^{\Gamma_0}=0\), and zero on \(H\); hence it is zero
on \(G=R_0H\).
Characters of \(G\) vanishing on \(H\) are the continuous
characters of \(B\). The criterion and uniqueness therefore
give \eqref{eq:all-cm-central-kernel-identification}.
\end{proof}

For the collection of all prime degrees, this theorem
eliminates the conjugation obstruction in
Proposition~\ref{prop:all-cm-intrinsic-extension}.
The relative central class remains a condition.
No arithmetic factorization \(G=JH\), equality of the
individual full and arithmetic scalar images, or
vanishing of \(\operatorname{Hom}_{\rm cont}(B,\mathbb F_2)\)
follows from these results.

